By Appendix~C (Frobenius orbits), as $\chi$ runs over the characters of $G_n$ with $\chi(\tau)=-1$, i.e.\ the even characters of conductor exactly $q=2^{n+2}$, the value $\chi(3)$ runs exactly once over the primitive $2^n$-th roots of unity $\omega$: indeed $\chi$ is determined by $\chi(3)$ because $3$ generates $G_n$, and $\chi(3)^m=\chi(\tau)=-1$ with $m=2^{n-1}$ says that $\chi(3)$ is a root of $x^m+1=\Phi_{2^n}(x)$, all of whose $m$ roots are primitive $2^n$-th roots of unity. Hence $\prod_\chi(1-\chi(3))=\prod_{\omega\text{ primitive}}(1-\omega)=\Phi_{2^n}(1)=1^{m}+1=2$, and taking absolute values, $\prod_\chi|1-\chi(3)|=2$ for every $n\ge2$. (The symbolic check in $\mathbb Q(\zeta_{2^n})$ for $n\le12$ in \path{sage/r11_propD_audit.log} is a numerical confirmation of this identity, not its proof.)
